\section{Related Work}
\label{sec:related}

\subsection{IoU-family segmentation metrics}
The region-overlap baseline this paper argues against for linear features. Region-overlap IoU is
the incumbent evaluation convention across general-purpose semantic-segmentation benchmarks ---
PASCAL VOC \cite{everingham2010pascalvoc} and Cityscapes \cite{cordts2016cityscapes} --- neither of
which contains a thin, 1-pixel-wide linear class, so neither benchmark's own leaderboard practice
had to confront IoU's road-width sensitivity at all. \textbf{Motivational baseline (see also
Section~\ref{subsec:related-roads}):} on the ISPRS Potsdam benchmark
\cite{rottensteiner2012isprs,rottensteiner2014isprs}, which \emph{does} include a building class
close in kind to PLEM's polygon classes, modern models report mean IoU up to ${\sim}86\%$ and mean
F1 up to ${\sim}92\%$ --- i.e.\ plain region-overlap metrics are already near-saturated on
well-behaved polygon classes. This is the same ``good-looking score, hidden failure mode'' pattern
Section~\ref{subsec:related-roads}'s road numbers show more starkly, and is why this paper does not
treat a high IoU/F1 figure alone as evidence a metric is adequate for PLEM's target feature types.

\subsection{Topology-/skeleton-aware metrics for linear structures}
clDice and related centerline methods (Shit et~al., CVPR 2021 \cite{shit2021cldice}).

\subsection{Boundary-based metrics for polygonal features}
Boundary F1 / BF score (Csurka et~al., BMVC 2013 \cite{csurka2013evaluation}).

\subsection{Road-network extraction and connectivity metrics}
\label{subsec:related-roads}
APLS (the SpaceNet road-extraction challenge metric \cite{vanetten2018spacenet}) and TOPO
(Biagioni \& Eriksson \cite{biagioni2012topo}), the two dominant graph-connectivity metrics for
linear road networks, motivate Sections~\ref{subsec:apls} and~\ref{sec:topological}'s
\texttt{apls}/\texttt{dtaf1\_topo} primitives. \textbf{Motivational baseline:} published
road-extraction leaderboards plateau well short of 1.0 even at their best --- SpaceNet Challenge
3's winning submission scored APLS 0.666 (field topped out ${\approx}0.67$); CRESI, the
best-known raster-to-graph pipeline, reaches APLS $0.69 \pm 0.02$ \cite{vanetten2019cresi};
D-LinkNet, winner of the DeepGlobe 2018 Road Extraction Challenge, reports pixel-IoU 0.647 (val) /
0.634 (test) \cite{zhou2018dlinknet} --- a \emph{pixel}-only ceiling that says nothing about
whether the extracted network stays connected. SpaceNet's building-footprint challenges show the
same pattern from the other direction: the first SpaceNet building challenge topped out at F1
0.26, and the four-city SpaceNet~2 challenge's winner reached F1 0.693 averaged across cities
(both per \cite{vanetten2018spacenet} and the associated challenge results reporting). These are
cited here as \textbf{motivating context, not as numbers PLEM's own trained models are benchmarked
against} --- PLEM's joint-training notebooks evaluate a small curated tile subset under this
paper's own protocol, not the official SpaceNet/DeepGlobe/Potsdam test sets, so the two sets of
numbers are not directly comparable. The point of citing them is narrower and stronger than a
leaderboard comparison: even the \emph{best} published pixel/topology scores on real road and
building extraction plateau in the 0.26--0.93 range depending on class and difficulty, and a
single scalar in that range cannot by itself distinguish scattered pixel noise within tolerance
from a genuinely fragmented network --- exactly the gap \texttt{dtaf1\_topo}/\texttt{apls} close
(see the road-breakage finding in Section~\ref{sec:topological}, where \texttt{dtaf1} stays pinned
at 1.0 under 75\% real road-pixel deletion while \texttt{dtaf1\_topo}/\texttt{apls} collapse
sharply).

\subsection{Point/instance detection metrics}
COCO-style average precision \cite{lin2014coco}, and the optimal bipartite-matching literature
underlying \texttt{point\_f1}'s Hungarian assignment. Citation target: Kuhn, 1955 (the Hungarian
algorithm) \cite{kuhn1955hungarian}.

\subsection{Composite/multi-task segmentation evaluation}
Positions the gap PLEM's \emph{evaluation} half fills: existing work evaluates linear and
polygonal feature types with separate, non-comparable metrics rather than a single class-agnostic
formula (DTAF1) plus a deliberately harsh cross-type composite (CBHM).

\subsection{Differentiable segmentation losses}
\label{subsec:related-losses}
The training-side analog of the gap the previous subsection describes on the eval side. Standard
multiclass segmentation losses (cross-entropy, soft Dice --- citation target: Milletari et~al.,
3DV 2016, ``V-Net'' \cite{milletari2016vnet}) are geometry-agnostic in exactly the sense
Section~\ref{sec:intro} argues IoU is: no notion of positional tolerance, centerline topology, or
point-instance identity. Three prior lines of work each address one of those properties in
isolation, none jointly:
\begin{enumerate}
  \item \textbf{Boundary-/Hausdorff-distance losses}, which precompute a GT-side distance
    transform and use it as a fixed per-pixel weight against the model's soft output --- citation
    targets: Kervadec et~al., 2019 (``Boundary loss for highly unbalanced segmentation'')
    \cite{kervadec2019boundary}; Karimi \& Salcudean, 2019 (Hausdorff-distance loss)
    \cite{karimi2019hausdorff} --- the direct precedent for Section~\ref{subsec:tolerance-loss}'s
    soft tolerance-band term.
  \item \textbf{Soft-clDice}, a differentiable soft-skeletonization loss for tubular structures ---
    citation target: Shit et~al., CVPR 2021 \cite{shit2021cldice} (the same paper cited for clDice
    itself in Section~\ref{subsec:cldice} --- here cited for its \emph{loss}, not its eval metric,
    contribution) --- the direct precedent for Section~\ref{subsec:cldice-loss}.
  \item \textbf{Heatmap regression for keypoint/point detection}, which reformulates discrete
    point detection as dense Gaussian-target regression rather than an instance-matching problem
    --- citation targets: Law \& Deng, ECCV 2018 (``CornerNet'') \cite{law2018cornernet};
    Zhou et~al., 2019 (``CenterNet''/``Objects as Points'') \cite{zhou2019centernet} --- the
    direct precedent for Section~\ref{subsec:heatmap-loss}.
\end{enumerate}
Section~\ref{sec:method}'s contribution is combining well-precedented individual relaxations of
(1)--(3) into one multi-task loss covering all three geometric properties simultaneously, plus the
masked multi-source training scheme (Section~\ref{subsec:multitask-loss}) that lets two datasets
which individually annotate only two of the three feature types jointly supervise one model.
